\documentclass[12pt]{article}
\usepackage[utf8]{inputenc}
\usepackage[T1]{fontenc}
\usepackage[margin=1in]{geometry}
\usepackage{setspace}
\usepackage{xcolor}
\usepackage[colorlinks=true, urlcolor=blue!60!black, linkcolor=blue!60!black]{hyperref}
\setlength{\parindent}{2em}
\setlength{\parskip}{1.25em}
\renewcommand{\baselinestretch}{1.0}
\begin{document}
\begin{center}
{\LARGE \textit{Citizens United v. Federal Election Commission}, 558 U.S. 310 (2010)\par}
\vspace{0.8em}
{\large Case Brief\par}
\end{center}
\vspace{1em}

\section{Memory Jogger}

In \href{https://www.courtlistener.com/opinion/1741/citizens-united-v-federal-election-commission/}{\textit{Citizens United v. Federal Election Commission}}, a 5--4 Supreme Court struck down federal bans on corporate and union independent political expenditures under the First Amendment while upholding campaign finance disclaimer and disclosure requirements by an 8--1 vote.

\section{Detailed Case Facts}

Citizens United, a conservative nonprofit corporation with an annual budget of roughly \$12 million funded mostly by individual donors and in small part by for-profit corporations, released a 90-minute documentary titled \textit{Hillary: The Movie} in January 2008. The film sharply criticized then-Senator Hillary Clinton, a leading candidate in the 2008 Democratic presidential primaries, featuring interviews with political commentators questioning her character and fitness for the presidency. Although Citizens United distributed the film in theaters and on DVD, it sought to expand reach by paying a cable provider \$1.2 million to make \textit{Hillary} available free of charge on a digital cable video-on-demand channel (``Elections '08'') reaching 34.5 million subscribers nationwide, and by airing two 10-second and one 30-second broadcast and cable television advertisements promoting the film.

Citizens United feared civil and criminal liability under 2 U.S.C. \S~441b, as amended by \S~203 of the Bipartisan Campaign Reform Act of 2002 (BCRA), commonly known as the McCain--Feingold Act. Section 441b prohibited corporations and labor unions from using general treasury funds for ``express advocacy'' or for any ``electioneering communication''---defined under 2 U.S.C. \S~434(f)(3)(A) as a broadcast, cable, or satellite communication referring to a clearly identified federal candidate within 30 days of a primary election or 60 days of a general election and capable of reaching 50,000 or more persons in the relevant electorate under 11 C.F.R. \S~100.29. Corporations and unions remained permitted to establish a ``separate segregated fund'' (a political action committee, or PAC) funded by restricted donations from stockholders, employees, or union members under 2 U.S.C. \S~441b(b)(2). Separately, BCRA \S\S~201 and 311 (2 U.S.C. \S\S~434(f) and 441d) mandated televised disclaimers identifying the sponsor and stating that the communication was not candidate-authorized, along with FEC disclosure filings for expenditures exceeding \$10,000 in a calendar year.

Historically, federal restrictions on corporate political spending dated back to the Tillman Act of 1907 (barring direct corporate contributions) and the Taft--Hartley Act of 1947 (extending restrictions to independent expenditures). In \textit{Buckley v. Valeo}, 424 U.S. 1 (1976), the Court upheld limits on direct candidate contributions to prevent \textit{quid pro quo} corruption but struck down ceilings on independent expenditures; two years later, \textit{First National Bank of Boston v. Bellotti}, 435 U.S. 765 (1978), invalidated a state ban on corporate spending in ballot referenda. In \textit{Austin v. Michigan Chamber of Commerce}, 494 U.S. 652 (1990), however, the Court upheld a state ban on corporate independent expenditures in candidate elections based on an ``antidistortion'' interest in curbing immense aggregations of corporate wealth. The Court subsequently relied on \textit{Austin} in \textit{McConnell v. Federal Election Commission}, 540 U.S. 93 (2003), to uphold BCRA \S~203 on its face, before narrowing \S~203 in \textit{FEC v. Wisconsin Right to Life, Inc.} (\textit{WRTL}), 551 U.S. 449 (2007), to cover only express advocacy or its ``functional equivalent.'' See \href{https://supreme.justia.com/cases/federal/us/558/310/}{Justia Opinion Summary} and \href{https://www.fec.gov/legal-resources/court-cases/citizens-united-v-fec/}{FEC Case Overview}.

\section{Procedural History}

In December 2007, Citizens United sued the Federal Election Commission in the U.S. District Court for the District of Columbia seeking declaratory and injunctive relief, arguing that 2 U.S.C. \S~441b was unconstitutional as applied to \textit{Hillary} and that BCRA \S\S~201 and 311 were unconstitutional as applied to the film and its promotional advertisements. Citizens United initially included a facial challenge to \S~441b (Count 5) but later stipulated to dismissing that count.

A three-judge District Court panel (Circuit Judge A. Raymond Randolph and District Judges Royce C. Lamberth and Richard W. Roberts) denied a preliminary injunction in \textit{Citizens United v. FEC}, 530 F. Supp. 2d 274 (D.D.C. 2008) (per curiam), holding that \S~441b was facially valid under \textit{McConnell}, that \textit{Hillary} was the functional equivalent of express advocacy under \textit{WRTL} because it was susceptible of no reasonable interpretation other than urging viewers to vote against Senator Clinton, and that BCRA's disclaimer and disclosure provisions were valid under \textit{McConnell}. On July 18, 2008, the three-judge panel granted summary judgment to the FEC on the same reasoning.

Pursuant to BCRA's direct-appeal provision, Citizens United appealed directly to the U.S. Supreme Court, which noted probable jurisdiction on November 14, 2008 (\href{https://www.oyez.org/cases/2008/08-205}{Oyez Case Record}). Following initial oral argument on March 24, 2009, the Court ordered supplemental briefing and reargument (held September 9, 2009) on whether to overrule \textit{Austin} and the portion of \textit{McConnell} upholding BCRA \S~203's facial validity. On January 21, 2010, the Supreme Court reversed in part, affirmed in part, and remanded.

\section{Judicial Votes}

\begin{itemize}
  \item \textbf{Majority (5--4 on Independent Expenditure Ban; 8--1 on Disclaimer and Disclosure Provisions):} Justice Kennedy delivered the opinion of the Court, joined in full by Chief Justice Roberts and Justices Scalia and Alito. Justice Thomas joined all parts of the opinion except Part IV (making Parts I--III and V a 5--4 majority). Justices Stevens, Ginsburg, Breyer, and Sotomayor joined Part IV (making Part IV an 8--1 majority upholding BCRA \S\S~201 and 311).
  \item \textbf{Concurrences:}
  \begin{itemize}
    \item Chief Justice Roberts filed a concurring opinion, joined by Justice Alito.
    \item Justice Scalia filed a concurring opinion, joined by Justice Alito, and joined in part by Justice Thomas.
  \end{itemize}
  \item \textbf{Opinions Concurring in Part and Dissenting in Part:}
  \begin{itemize}
    \item Justice Stevens filed an opinion concurring in part (as to Part IV) and dissenting in part (as to Parts I--III and V), joined by Justices Ginsburg, Breyer, and Sotomayor.
    \item Justice Thomas filed an opinion concurring in part (as to Parts I--III and V) and dissenting in part (as to Part IV).
  \end{itemize}
\end{itemize}

\section{Holding}

\textbf{Judgment:} Reversed in part as to the constitutionality of 2 U.S.C. \S~441b's prohibition on corporate independent expenditures; affirmed in part as to the constitutionality of BCRA \S\S~201 and 311's disclaimer and disclosure requirements as applied to \textit{Hillary} and its promotional ads; and remanded.

\textbf{Rules Announced:}
\begin{enumerate}
  \item Under the First Amendment, the government may not suppress political speech on the basis of the speaker's corporate identity; laws banning independent political expenditures by corporations (and, by extension, labor unions) are subject to strict scrutiny and fail because independent expenditures do not give rise to \textit{quid pro quo} corruption or the appearance of corruption. \textit{Austin v. Michigan Chamber of Commerce}, 494 U.S. 652 (1990), and the portion of \textit{McConnell v. FEC}, 540 U.S. 93 (2003), upholding BCRA \S~203 are overruled.
  \item Campaign finance disclaimer and disclosure requirements are subject to ``exacting scrutiny'' (requiring a substantial relation to a sufficiently important governmental interest) and are constitutional as applied to electioneering communications and their promotional ads based on the government's informational interest, absent a showing of a reasonable probability of threats, harassment, or reprisals against donors.
\end{enumerate}

\textbf{Holding vs. Dicta:} The Court's invalidation of 2 U.S.C. \S~441b's corporate independent expenditure ban, its overruling of \textit{Austin} and \textit{McConnell}'s \S~203 holding, and its affirmance of BCRA \S\S~201 and 311 as applied to Citizens United are binding holdings. By contrast, the Court expressly declined to decide whether the government has a compelling interest in restricting political spending by foreign nationals or foreign-controlled corporations (noting \S~441b was overbroad regardless), and it did not disturb existing limits on direct contributions to candidates, which Citizens United did not challenge.

\section{Analysis of Opinions}

\subsection{Majority Opinion (Justice Kennedy)}

Writing for the Court (\href{https://www.law.cornell.edu/supct/html/08-205.ZO.html}{Cornell LII Full Text}), Justice Kennedy first rejected Citizens United's narrower statutory and as-applied arguments. Under 11 C.F.R. \S~100.29, \textit{Hillary} was ``publicly distributed'' because the cable video-on-demand system reached 34.5 million subscribers; under \textit{WRTL}, the film was the ``functional equivalent of express advocacy'' because it was a feature-length negative appeal to vote against Senator Clinton; courts could not draw durable constitutional lines around video-on-demand technology; and creating a \textit{de minimis} exception for nonprofits accepting small corporate donations would contradict the text of \textit{FEC v. Massachusetts Citizens for Life, Inc.} (\textit{MCFL}), 479 U.S. 238 (1986), and BCRA's Snowe--Jeffords Amendment. Because Citizens United preserved a First Amendment claim and the District Court passed upon facial validity, the Court held that it had to confront \textit{Austin} rather than prolong a complex FEC regulatory regime that operated as a \textit{de facto} prior restraint chilling core political speech.

On the merits, Kennedy treated \S~441b as an outright ban on corporate speech backed by felony penalties, rejecting the claim that the PAC option cured the defect because a PAC is a separate, heavily regulated entity. Applying strict scrutiny, the Court held that the First Amendment prohibits speaker-based restrictions on political speech and dismantled all three governmental justifications for \S~441b:
\begin{itemize}
  \item \textbf{Antidistortion Rationale:} Kennedy rejected \textit{Austin}'s theory that government may curb ``immense aggregations of wealth'' accumulated through the corporate form. Under \textit{Buckley} and \textit{Davis v. FEC}, 554 U.S. 724 (2008), the government has no legitimate interest in equalizing speakers' voices. Moreover, \textit{Austin}'s logic would permit Congress to censor media corporations---which are statutorily exempt only by legislative grace---and swept in millions of small and nonprofit corporations.
  \item \textbf{Anticorruption Rationale:} Relying on \textit{Buckley}, Kennedy confined the compelling anticorruption interest strictly to \textit{quid pro quo} corruption (dollars for political favors), declaring that ``independent expenditures, including those made by corporations, do not give rise to corruption or the appearance of corruption.'' Mere ingratiation, influence, or access is an inherent feature of representative democracy, not corruption, and \textit{Caperton v. A. T. Massey Coal Co.}, 556 U.S. 868 (2009), was distinguished as a Due Process judicial-recusal case rather than a speech-suppression case.
  \item \textbf{Shareholder Protection:} The Court found the shareholder-protection rationale both underinclusive (covering only broadcast/cable/satellite ads within 30 or 60 days of an election) and overinclusive (reaching nonprofits and single-shareholder corporations), concluding that dissenting shareholders can rely on ``the procedures of corporate democracy.''
\end{itemize}

Applying \textit{stare decisis} factors, Kennedy concluded that \textit{Austin} was poorly reasoned, abandoned even by the government's own defense, undermined by technological and campaign-finance realities, and supported by no serious reliance interests. Finally, in Part IV (joined by eight Justices), Kennedy upheld BCRA \S\S~201 and 311 under exacting scrutiny, emphasizing that disclosure imposes no ceiling on speech, informs voters and shareholders in the Internet age, and leaves open as-applied exemptions where donors face a reasonable probability of threats or harassment.

\subsection{Concurring Opinion (Chief Justice Roberts)}

Joined by Justice Alito, Chief Justice Roberts wrote separately to defend the Court's adherence to judicial restraint and \textit{stare decisis}. Roberts argued that constitutional avoidance requires adopting a narrower ground only when it is legally sound: ``There is a difference between judicial restraint and judicial abdication.'' Once Citizens United's narrower statutory and as-applied arguments failed on the merits, the Court had no choice but to address \textit{Austin}. On \textit{stare decisis}, Roberts emphasized that \textit{Austin} itself was an aberration departing from \textit{Buckley} and \textit{Bellotti}, had generated persistent doctrinal instability, threatened press freedoms if carried to its logical conclusion, and could not be preserved by ``jury-rigging new and different justifications'' that the \textit{Austin} Court itself never adopted.

\subsection{Concurring Opinion (Justice Scalia)}

Joined by Justice Alito and in part by Justice Thomas, Justice Scalia directly rebutted Justice Stevens's originalist critique. Scalia argued that the First Amendment's text protects ``speech,'' not specific categories of ``speakers,'' and contains no exception for associations organized in corporate form. He contended that roughly 335 business corporations were chartered in the United States by 1800 alongside municipal, religious, and educational corporations that routinely petitioned and published, and that founding-era hostility was aimed at state-granted monopolies rather than modern general-incorporation entities. Scalia also argued that the Free Press Clause historically protected the publishing technology available to all citizens rather than a special privilege for institutional media, making the dissent's distinction between media and non-media corporations textually unsustainable.

\subsection{Opinion Concurring in Part and Dissenting in Part (Justice Stevens)}

Joined by Justices Ginsburg, Breyer, and Sotomayor, Justice Stevens concurred in Part IV upholding disclosure but delivered a sweeping 90-page dissent from the invalidation of \S~441b. Procedurally, Stevens accused the majority of abandoning judicial restraint by reviving a facial challenge that Citizens United had expressly dismissed, ordering reargument to manufacture a vehicle to overrule \textit{Austin} and \textit{McConnell} without an evidentiary record, and discarding narrower options (such as a video-on-demand carve-out or an expanded \textit{MCFL} exemption).

Substantively, Stevens argued that \S~441b was not a ``ban'' on speech because corporations could speak freely outside the 30/60-day window or at any time through a PAC. He rejected the majority's ``glittering generality'' against speaker-based distinctions, noting that First Amendment doctrine regularly accounts for speaker identity (e.g., students, prisoners, military personnel, and foreign nationals). Structurally, corporations are artificial legal entities with limited liability, perpetual life, and fiduciary duties to maximize profit; they cannot vote, run for office, or belong to ``We the People,'' and they may be controlled by nonresidents or foreign shareholders. Drawing on founding-era history, Stevens argued that the Framers viewed corporations as tightly circumscribed creatures of state charter and understood ``corruption'' far more broadly than \textit{quid pro quo} bribery to include undue corporate domination of the legislative process and erosion of public trust. He warned that the majority's ruling overturned a century of federal and state campaign finance law and would unleash corporate treasury spending that threatens the integrity of American democracy.

\subsection{Opinion Concurring in Part and Dissenting in Part (Justice Thomas)}

Justice Thomas joined the majority in striking down \S~441b but dissented from Part IV's affirmance of BCRA \S\S~201 and 311 (\href{https://www.law.cornell.edu/supct/html/08-205.ZX1.html}{Cornell LII Opinion}). Grounding his argument in the First Amendment's historical protection of anonymous political speech (\textit{McIntyre v. Ohio Elections Commission}, 514 U.S. 334 (1995)), Thomas documented extensive real-world retaliation against disclosed donors to California's Proposition 8---including online maps targeting donors' homes and businesses, death threats, boycotts, forced resignations, and property damage---as well as threats of retaliation by elected officials. Thomas argued that it was ironic for the majority to reject case-by-case as-applied challenges as too chilling under \S~203 while insisting that those same slow as-applied challenges adequately protect donors from intimidation under \S\S~201 and 311 in the Internet era.

\section{Examples: Future Applications}

\begin{enumerate}
  \item \textbf{Same-Side (Invalidated): State Ban on Corporate Independent Expenditures.} A state enacts a statute prohibiting for-profit corporations from spending general treasury funds on independent television advertisements supporting or opposing candidates for state legislature. Under \textit{Citizens United}, the statute is facially unconstitutional because the First Amendment bars states---just as it bars Congress---from suppressing independent political expenditures based on the speaker's corporate identity. (The Supreme Court later applied this exact principle summarily in \textit{American Tradition Partnership, Inc. v. Bullock}, 567 U.S. 516 (2012).)
  \item \textbf{Same-Side (Invalidated): Contribution Caps on Independent-Expenditure-Only Committees (Super PACs).} Federal law caps individual contributions to a political committee that makes solely independent expenditures and never coordinates with or contributes directly to candidates. Under \textit{Citizens United}'s categorical holding that independent expenditures do not give rise to \textit{quid pro quo} corruption or its appearance, contributions to an entity that engages exclusively in independent spending likewise lack a corrupting nexus, rendering the cap unconstitutional. (This reasoning directly drove \textit{SpeechNow.org v. FEC}, 599 F.3d 686 (D.C. Cir. 2010).)
  \item \textbf{Opposite-Side (Upheld): Federal Ban on Direct Corporate Contributions to Candidates.} A business corporation challenges 2 U.S.C. \S~441b's remaining ban on using general treasury funds to make direct monetary contributions to a congressional candidate's campaign committee. Under \textit{Citizens United}, the ban remains constitutional because the Court preserved \textit{Buckley}'s distinction between independent expenditures and direct contributions, recognizing that direct contributions pose a concrete risk of \textit{quid pro quo} corruption and circumvention of individual contribution limits.
  \item \textbf{Opposite-Side (Upheld): Mandatory Sponsorship Disclaimers on Corporate Digital Election Ads.} A state requires any corporation spending more than \$5,000 on online video advertisements advocating for or against a state ballot candidate within 60 days of an election to display a ``Paid for by [Organization]'' disclaimer and file donor reports with the state election board. Under Part IV of \textit{Citizens United}, the law is constitutional under exacting scrutiny because it imposes no ceiling on spending and advances the sufficiently important informational interest of helping voters evaluate political messages, unless the group proves a reasonable probability of donor harassment.
  \item \textbf{Fence-Sitter (Close Call): Ban on Independent Political Spending by Foreign-Owned Domestic Subsidiaries.} Congress enacts a statute prohibiting any U.S.-incorporated subsidiary whose shares are more than 50\% owned by foreign nationals or foreign governments from making independent expenditures in federal elections. \textit{Citizens United} expressly left open whether the government has a compelling interest in preventing foreign influence over American elections (Part III-B-4), creating a sharp tension between the Court's anti-speaker-discrimination principle (since a U.S. subsidiary is a domestic corporate entity employing American citizens) and the sovereignty-based interest in excluding foreign nationals from democratic self-government (later addressed as to direct foreign nationals in \textit{Bluman v. FEC}, 800 F. Supp. 2d 281 (D.D.C. 2011), \textit{aff'd}, 565 U.S. 1104 (2012)).
\end{enumerate}

\section{Critique}

\textit{Citizens United} is among the most fiercely debated constitutional decisions of the twenty-first century, drawing criticism from both progressive and originalist or textualist scholars, alongside vigorous defenses from First Amendment libertarians.

From a \textbf{progressive and democratic-theory perspective}, scholars such as Richard L. Hasen (\href{https://repository.law.umich.edu/context/mlr/article/1166/viewcontent/Hasen_CITIZENS_20UNITED_20AND_20THE_20ILLUSION_20OF_20COHERENCE.pdf}{\textit{Citizens United and the Illusion of Coherence}}, 109 Mich. L. Rev. 581 (2011)) and organizations like the \href{https://www.brennancenter.org/our-work/research-reports/citizens-united-explained}{Brennan Center for Justice} argue that the majority's categorical empirical claim---that independent expenditures ``do not give rise to corruption or the appearance of corruption''---is flatly contradicted by political reality. Critics point out that the rise of single-candidate ``Super PACs'' funded by massive corporate and billionaire contributions blurs any functional line between ``independent'' and coordinated spending, enabling donors to secure outsized legislative access and ``dependency corruption.'' Progressive critics also fault the Court's reliance on ``corporate democracy'' and prompt disclosure as illusory safeguards, noting that subsequent growth in 501(c)(4) ``dark money'' groups obscured true funding sources and that institutional investors rarely check management's political spending.

From an \textbf{originalist and historical-textualist perspective}, critics challenge the majority's claim that the 1791 First Amendment barred regulation of corporate electioneering. Legal historian Zephyr Teachout (\href{https://ir.lawnet.fordham.edu/cgi/viewcontent.cgi?article=1430&context=faculty_scholarship}{\textit{Historical Roots of Citizens United vs. FEC}}, 5 Harv. L. \& Pol'y Rev. 163 (2011)) argues that the Framers held a thick, structural conception of ``corruption'' extending far beyond narrow \textit{quid pro quo} bribery to include improper dependence and private capture of public offices. Similarly, originalist-oriented scholarship in the \href{https://scholarship.law.nd.edu/cgi/viewcontent.cgi?article=4640&context=ndlr}{Notre Dame Law Review} contends that founding-era corporations were narrowly chartered artificial entities subject to plenary state regulation, making it ahistorical to equate modern multinational business corporations with individual citizens under the original public meaning of the First Amendment. Even some conservative judicial minimalists criticized the majority's procedural aggressiveness in ordering reargument to decide a facial challenge that the appellant had abandoned below.

Conversely, \textbf{free-speech defenders} such as Joel Gora (\href{https://readingroom.law.gsu.edu/cgi/viewcontent.cgi?article=2215&context=gsulr}{\textit{The Campaign Finance Debate After Citizens United}}, 27 Ga. St. U. L. Rev. 935 (2011)) praise the decision as a faithful vindication of core First Amendment principles. They argue that \S~441b criminalized political advocacy not only by Fortune 500 firms but by millions of small businesses and nonprofit advocacy groups (such as the ACLU, Sierra Club, and NRA), and that the government's concession at oral argument that \S~441b could theoretically ban books revealed the authoritarian reach of \textit{Austin}'s antidistortion rationale. At bottom, the clash in \textit{Citizens United} reflects a deep conflict between competing constitutional premises: a libertarian view of the First Amendment as an absolute check against government gatekeeping in the marketplace of ideas versus a republican view that democratic self-government requires legislatures to protect electoral equality and institutional integrity from concentrated economic power.

\section{Key Quotations}

\begin{quote}
``The Government may regulate corporate political speech through disclaimer and disclosure requirements, but it may not suppress that speech altogether.''
\end{quote}
\noindent --- Justice Kennedy, Opinion of the Court, 558 U.S. at 319.

\begin{quote}
``For the reasons explained above, we now conclude that independent expenditures, including those made by corporations, do not give rise to corruption or the appearance of corruption.''
\end{quote}
\noindent --- Justice Kennedy, Opinion of the Court, 558 U.S. at 357.

\begin{quote}
``We return to the principle established in \textit{Buckley} and \textit{Bellotti} that the Government may not suppress political speech on the basis of the speaker's corporate identity. No sufficient governmental interest justifies limits on the political speech of nonprofit or for-profit corporations.''
\end{quote}
\noindent --- Justice Kennedy, Opinion of the Court, 558 U.S. at 365.

\begin{quote}
``The Amendment is written in terms of `speech,' not speakers. Its text offers no foothold for excluding any category of speaker, from single individuals to partnerships of individuals, to unincorporated associations of individuals, to incorporated associations of individuals---and the dissent offers no evidence about the original meaning of the text to support any such exclusion.''
\end{quote}
\noindent --- Justice Scalia, concurring, 558 U.S. at 392--93.

\begin{quote}
``In the context of election to public office, the distinction between corporate and human speakers is significant. Although they make enormous contributions to our society, corporations are not actually members of it. They cannot vote or run for office.''
\end{quote}
\noindent --- Justice Stevens, concurring in part and dissenting in part, 558 U.S. at 394.


\end{document}
